\documentclass[12pt]{article}
\usepackage[T1]{fontenc}
\usepackage[utf8]{inputenc}
\usepackage{lmodern}
\usepackage{textcomp}
\usepackage{enumitem}
\usepackage{geometry}
\usepackage{graphicx}
\usepackage{amsmath, amssymb}
\geometry{margin=1in}
\begin{document}
\begin{center}
Physics - Graduate Level Assessment 11 \
Total Marks: 40 \
Answer all questions.
\end{center}

\section*{Multiple Choice (10 marks)}
\begin{enumerate}[label=\arabic*.]
\item A principal source of the Earth's internal energy is
\begin{enumerate}[label=(\alph*)]
    \item electromagnetic radiation.
    \item thermal conduction.
    \item gravitational pressure.
    \item radioactivity.
    \item nuclear fusion.
\end{enumerate}
\item The negative muon, mu^-, has properties most similar to which of the following?
\begin{enumerate}[label=(\alph*)]
    \item Gluon
    \item Electron
    \item Quark
    \item Proton
    \item Photon
\end{enumerate}
\item How many nuclei are there in 1 kg of aluminum?
\begin{enumerate}[label=(\alph*)]
    \item 7.21 x $10^25$ nuclei
    \item 9.09 x 10^-26 nuclei
    \item 4.48 x 10^-26 nuclei
    \item 26.98 x 10^-27 nuclei
    \item 3.01 x $10^22$ nuclei
\end{enumerate}
\item The force required to start a mass of 50 kilograms moving over a rough surface is 343 Nt. What is the coefficient of starting friction?
\begin{enumerate}[label=(\alph*)]
    \item 0.75
    \item 0.40
    \item 0.55
    \item 0.65
    \item 0.50
\end{enumerate}
\item How are planetary rings made?
\begin{enumerate}[label=(\alph*)]
    \item From accretion within the solar nebula at the same time the planets formed
    \item From fragments of planets ejected by impacts
    \item From dust grains that escape from passing comets
    \item From the remnants of a planet's atmosphere
    \item From the gravitational pull of nearby planets
\end{enumerate}
\end{enumerate}

\section*{True / False (10 marks)}
\textit{Answer True or False}
\begin{enumerate}[label=\arabic*.]
\setcounter{enumi}{5}
\item True or False: The centripetal force needed for the car to turn is significantly less than its weight, approximately 400 lb.
\item True or False: A minimum voltage of 65000 volts is sufficient to produce K radiation from a tungsten target in an x-ray tube.
\item True or False: The angular acceleration of the turntable, when accelerated from rest to 33.3 rpm in 2 seconds, is 1.74 sec^-2.
\item True or False: If the absolute temperature of a blackbody is increased by a factor of 3, the energy radiated per second per unit area increases by a factor of 81.
\item True or False: The magnetic field lines around a current-carrying wire form concentric circles that are oriented perpendicular to the wire.
\end{enumerate}

\section*{Long Form Response (20 marks)}
\textit{Answer each question with a clear and complete explanation.}
\begin{enumerate}[label=\arabic*.]
\setcounter{enumi}{10}
\item Discuss the relationship between voltage and energy dissipation in a resistor. What happens to the rate of energy dissipation if the voltage is doubled? Provide a detailed analysis.
\item Calculate the time period of each oscillation of an ultrasonic transducer operating at 6.7 MHz and discuss how to derive its angular frequency in rad/s.
\end{enumerate}

\vfill
\noindent\textit{End of Paper}
\end{document}